\section{The genuine sheaf of original module morphisms}

\begin{theorem}[Module-linearity descends along a covering sieve]
\label{thm:module-hom-original-linearity-descent}
\lean{EnvelopeProperty.ModuleHomSheaf.ambientHom,EnvelopeProperty.ModuleHomSheaf.toAmbient,EnvelopeProperty.ModuleHomSheaf.toAmbient_injective,EnvelopeProperty.ModuleHomSheaf.restrictHom,EnvelopeProperty.ModuleHomSheaf.ambientHom_map_toAmbient,EnvelopeProperty.ModuleHomSheaf.homSubfunctor,EnvelopeProperty.ModuleHomSheaf.homSubfunctor_isSheaf,EnvelopeProperty.ModuleHomSheaf.localHomSheaf,EnvelopeProperty.ModuleHomSheaf.localHomEquiv}
\leanok
Let $R$ be a sheaf of rings on a site and $M,N$ original sheaves of
$R$-modules. The local module morphisms $M|_U\to N|_U$ form a sheaf of
types. Its sections retain the entire original module morphisms, with
their original scalar actions.
\end{theorem}
\begin{proof}
Forget scalar multiplication and form the existing sheaf of morphisms
between the underlying additive sheaves. Inside it take the subpresheaf
of those local morphisms that are $R$-linear. Restriction preserves
linearity because both source and target restriction functors preserve
the original scalar actions.

It remains to show that this subpresheaf is closed under gluing. Suppose
an additive local morphism becomes linear along a covering sieve.
On any object over $U$, to compare its values on $rm$ and $r$ times its
value on $m$, pull the sieve back to that object. The target module
sheaf is separated for this covering sieve. On each covering arrow,
naturality moves restriction through the additive morphism. Both
source and target restriction maps preserve scalar multiplication, so
the desired equality becomes the scalar-linearity equality for one of
the original covering module morphisms. Separatedness proves it before
restriction. Thus the glued additive morphism is itself linear. A
subpresheaf of a sheaf closed under this local test is a sheaf, proving
the assertion.
\end{proof}
